Este documento consolida el desarrollo técnico de \textbf{GeoCebada} para el Reto
AgroCebada FIRA 2026. El problema operativo consiste en reconstruir
\artifact{RENDIMIENTO_T_HA} para un conjunto fijo de 59 parcelas, disponiendo de 138
parcelas etiquetadas y covariables observables para las 197. La formulación es
\emph{transductiva de objetivo fijo}: los $X$ de las 59 parcelas objetivo pueden intervenir
en transformaciones estrictamente $X$-only, mientras sus rendimientos permanecen ocultos.

Checkpoint~01 fijó la semántica y el contrato reproducible de las fuentes. Checkpoint~02
comprimió series satelitales y fuentes ambientales a una tabla parcelaria de 197 filas y
1,403 predictores. Checkpoint~03 añadió 350 variables agronómicas, 335 variables empíricas
$X$-only, descubrimiento fold-local y benchmarks anidados de representaciones, familias y
ensambles globales.

Checkpoint~04 reformuló la validación para reproducir la información real del reto. La señal
dominante resultó espacial: Moran's $I=0.6797$ con $p=0.001$ y Spearman $\rho=0.4882$ entre
distancia geográfica y diferencia absoluta de rendimiento. La búsqueda local/grafo produjo
\artifact{LocalRidge_C4_all_deterministic_Geo_k24_a30_p1} (Local04D), con RMSE
target-matched medio 0.487845 y pooled 0.495741.

Checkpoint~05 combinó globales y transductivos mediante cross-fitting, stacking y un
mixture-of-experts condicionado en soporte. El mejor mínimo de desarrollo, LocalE13 con
10\% de peso global, alcanzó 0.487243 de RMSE medio y 0.495261 pooled, pero el selector
leave-one-development-split-out obtuvo 0.487934 y 0.495933, ambos ligeramente peores que
Local04D. El gate falló y la confirmación fresca no se puntuó. La fuente canónica posterior
a 05 es \artifact{reports/checkpoint_05/final_predictions.csv} y conserva Local04D.

Después de ese cierre se ejecutó Checkpoint~03D como experimento aditivo de capacidad global.
La configuración final mantuvo los dos protocolos outer de cinco folds y tres folds inner,
pero redujo tuning redundante para evitar un run proyectado de varios días. El benchmark
balanced completó 230 outer fits en 248.085 minutos. Sus finalistas fueron HistGB, XGBoost y
LightGBM, todos sobre G1 agronomic, con RMSE municipality-grouped 0.715339, 0.733080 y
0.748478 respectivamente. El resultado mejora la frontera global grouped, pero no domina
simultáneamente state y grouped ni autoriza comparar directamente esos scores con el RMSE
target-matched de Local04D.

HistGB fue rank científico 1 pero no convertible con el stack ONNX estable. XGBoost rank 2
fue el finalista deployable de mayor rango; el round trip ONNX tuvo diferencia absoluta máxima
$3.814697\times10^{-6}$. Ese ONNX es un estimador global de deployment y no una codificación
exacta del predictor competition-facing Local04D.

La búsqueda global amplia queda cerrada. Cualquier último intento de combinar Local04D con
los nuevos globales 03D debe pasar por un Checkpoint~05B estrecho bajo el mismo protocolo
target-matched; de lo contrario, el trabajo siguiente es integración de submission, app y
deployment.
